% chapters/16/section-01.tex
\section{Mengapa Prototipe Harus Diuji}
\label{sec:mengapa-diuji}

\subsection{Pembuat Tidak Dapat Menilai Karyanya Sendiri}

Sepanjang Bab~III sampai Bab~XV, Anda menguji permainan Anda dengan satu cara
saja: menjalankannya, lalu memeriksa apakah perilakunya sesuai dengan yang Anda
maksudkan. Cara itu berhasil untuk menemukan kesalahan yang Anda sendiri sudah
tahu jawabannya, tetapi ia gagal pada hal yang paling penting.

Kegagalan itu terjadi karena Anda sudah tahu terlalu banyak. Anda tahu tombol
mana yang harus ditekan, tahu jalur mana yang harus dilewati, dan tahu ke mana
tujuan yang sebenarnya. Pengetahuan itu tidak dimiliki pemain baru, dan justru
pengetahuan itulah yang menutupi masalah. Sebuah level yang tampak jelas bagi
pembuatnya dapat membingungkan pemain lain selama beberapa menit pertama
\cite{unity-leveldesign}.

Karena itu, pengujian bukan pekerjaan yang dapat diselesaikan sendirian. Ia
menuntut orang lain yang belum pernah melihat karya Anda, dan menuntut Anda
menahan diri untuk tidak menolong mereka.

\begin{peringatanbox}
Kesalahan yang paling sering dilakukan mahasiswa pada sesi
\istilah{playtest} pertama: memberi petunjuk ketika pengujinya kebingungan.
Kalimat seperti ``tekan tombol spasi untuk lompat'' tampak seperti bantuan,
padahal ia menghapus satu-satunya keterangan yang paling Anda butuhkan.
Apabila penguji tidak menemukan sendiri caranya, itu bukan kesalahan penguji,
melainkan temuan.
\end{peringatanbox}

\subsection{Dua Jenis Pengujian yang Saling Melengkapi}

Ada dua jalan untuk memeriksa sebuah prototipe, dan keduanya menemukan hal yang
berbeda.

\begin{description}
  \item[\textbf{Pemeriksaan otomatis}] dijalankan oleh komputer, tanpa pemain,
        dan berulang-ulang. Ia menemukan hal yang bersifat pasti: petak yang
        tidak dapat dicapai, lompatan yang melebihi jangkauan, berkas yang
        gagal dimuat. Ia cepat, dan hasilnya sama setiap kali dijalankan.
  \item[\textbf{\istilah{Playtest}}] dijalankan oleh orang lain yang memainkan
        prototipe Anda. Ia menemukan hal yang tidak dapat dihitung: bagian yang
        membingungkan, bagian yang membosankan, dan bagian yang membuat pemain
        berhenti bermain.
\end{description}

Keduanya diperlukan, dan urutannya penting. Pemeriksaan otomatis dijalankan
lebih dahulu, sebab tidak ada gunanya meminta orang lain menguji level yang
belum dapat diselesaikan. Sekaligus, pemeriksaan otomatis tidak dapat
menggantikan \istilah{playtest}: mesin dapat memberi tahu Anda bahwa levelnya
\emph{mungkin} diselesaikan, tetapi tidak dapat memberi tahu apakah levelnya
\emph{menyenangkan} diselesaikan.

\begin{catatabox}
Pemeriksaan otomatis pada bab ini dijalankan dengan perkakas yang sudah Anda
miliki sendiri: Python, dan peta level yang Anda susun pada Bab~X. Tidak ada
perkakas tambahan yang perlu dipasang, dan hasilnya dapat diperiksa ulang.
\end{catatabox}

\subsection{Pengujian sebagai Siklus, Bukan sebagai Langkah Terakhir}

Kebiasaan yang perlu Anda tinggalkan adalah memandang pengujian sebagai
pekerjaan terakhir sebelum pengumpulan. Pengujian adalah siklus: setiap
perbaikan menuntut pemeriksaan ulang, dan setiap pemeriksaan ulang menuntut
daftar periksa yang diperbarui. Siklus itu diperlihatkan pada
Gambar~\ref{fig:siklus-uji}.

\begin{figure}[htbp]
\centering
% chapters/16/figures/siklus-uji.tikz
% Pengujian sebagai siklus: setiap perbaikan menuntut pengujian ulang.
\begin{tikzpicture}[
  tahap/.style={
    rectangle, rounded corners=5pt, draw=ubPrimary, thick, fill=ubAlt,
    align=center, text width=27mm, inner sep=4pt, minimum height=11mm,
    font=\footnotesize
  },
  alat/.style={
    rectangle, rounded corners=5pt, draw=ubAccent, thick, fill=white,
    align=center, text width=27mm, inner sep=4pt, minimum height=11mm,
    font=\footnotesize
  },
  panah/.style={-{Stealth[length=2.2mm]}, thick, draw=ubAccent},
  kembali/.style={-{Stealth[length=2.2mm]}, thick, draw=ubPrimary, dashed},
  catatan/.style={font=\scriptsize\itshape, align=left}
]

\node[tahap] (gdd)    at (0,0)     {\istilah{GDD}\\{\scriptsize\itshape rencana}};
\node[alat]  (daftar) at (3.9,0)   {checklist\\fungsional};
\node[alat]  (otomat) at (7.8,0)   {uji otomatis\\{\scriptsize\itshape tanpa pemain}};
\node[alat]  (main)   at (11.7,0)  {playtest\\{\scriptsize\itshape dengan pemain}};

\node[tahap] (log)    at (11.7,-2.0) {log \istilah{bug}};
\node[tahap] (perbaiki) at (7.8,-2.0) {perbaikan};
\node[tahap] (build)  at (3.9,-2.0) {build akhir};

\draw[panah] (gdd) -- (daftar);
\draw[panah] (daftar) -- (otomat);
\draw[panah] (otomat) -- (main);
\draw[panah] (main) -- (log);
\draw[panah] (log) -- (perbaiki);
\draw[panah] (perbaiki) -- (build);

% Kembali ke awal: perbaikan selalu menuntut pengujian ulang.
\draw[kembali] (perbaiki.north) .. controls (7.8,1.9) and (5.6,1.9) .. (daftar.north);

\node[catatan, anchor=north west, text width=62mm] at (-1.6,-2.9)
  {Perbaikan tidak mengakhiri pekerjaan, melainkan memulai putaran
   berikutnya. Garis putus-putus menandai bahwa setiap perbaikan menuntut
   pemeriksaan ulang, dan pemeriksaan ulang menuntut checklist yang
   diperbarui.};
\end{tikzpicture}

\caption{Pengujian sebagai siklus: setiap perbaikan menuntut pemeriksaan ulang}
\label{fig:siklus-uji}
\end{figure}

Gambar itu tidak berakhir pada \istilah{build} akhir,
melainkan kembali ke daftar periksa. Alasannya sederhana: perbaikan yang Anda
lakukan atas satu temuan berpotensi menimbulkan temuan baru. Menambahkan
pijakan agar pemain tidak terjebak dapat membuat bagian lain menjadi terlalu
mudah; menggeser letak tujuan dapat membuat satu bagian level tidak lagi
terjangkau. Tanpa pemeriksaan ulang, perbaikan justru menambah masalah.
